% $Id: chen.tex,v 1.2 2001/10/29 17:38:27 chenli Exp $

\subsection{Heterogeneous data integration}

\subsubsection*{Task 1: Extracting semi-structured data from web pages}

To support cross-modality search, we need to extract structured
information from web pages.  This extraction is also critical to many
other applications.  For instance, comparison-shopping engines need to
retrieve product information from their underlying sources.  In addition,
extracting web information effectively can meet the growing demand in
making web content available to mobile devices (e.g., PDAs and cell
phones), which have limited display and network bandwidth.

Many approaches have been proposed on extracting information from web
pages.  Many applications need software components called {\em scrapers}
to do the extracting automatically.  We see that the data is going through
a full cycle: data is stored in structured databases at publisher sites,
published in unstructured formats, and scraped and stored in structured
databases again at the consumer end.  Scrapers are typically built using
rule-based techniques.  Often, these techniques involve a human spending a
lot of time on the sites, figuring out the format and layout of the pages,
and then writing custom rules for scraping the data.  Apart from being
time consuming and error prone, these techniques are not robust to changes
in the webpage structure.

In our study, we plan to tackle this problem by analyzing the regularities
shared by different pages of the same web site.  For instance,
Figure~\ref{fig:bookpages} shows two book pages at Amazon.com.  The figure
clearly indicates the similarities of the two pages.  They share the same
layout, in which some portions have same contents.  One important
observation is that some interesting information, such as book title,
authors, price, and availability, are included in the differences.

\begin{figure}[htfb]
\begin{center}
\begin{tabular}[htfb]{ll}
\epsfxsize=200pt
\epsfysize=150pt
\subfigure[Book 1.]{\epsfbox{amazon1.eps}}
\epsfxsize=200pt
\epsfysize=150pt
\subfigure[Book 2.]{\epsfbox{amazon2.eps}}
\end{tabular}
\end{center}
\caption{Two book pages at Amazon.com.}
\label{fig:bookpages}
\end{figure}

In general, a lot of Web information is stored in the form of highly
structured, often dynamically generated pages.  These pages are typically
generated by plugging in data from a structured backend source (relational
database) into HTML templates.  Our approach is to infer the structure of
these HTML templates used to generate the webpages by examining many
webpages.  By analyzing these templates we can extract useful information
effectively.  Our approach works in the following steps:

\begin{enumerate}
  \item Choose some representative pages from a web site that looks
  similar.

  \item Extract the common template shared by these pages.

  \item Show the differences to users.  Let users choose the values that
  are useful.

  \item Modify the template.  Apply the new template to other pages to
  extract useful information to populate the result database.
\end{enumerate}

Our approach will work in the case where more than one template is used to
generate the pages.

\subsubsection*{Task 2: Constructing Meta-data for Multimedia Objects}

It is known that content-based search on multimedia information (e.g.,
images, audio, and video) has limitations, since it is hard to extract
features that can match concepts understandable by human being.  Hence
constructing meta data for multimedia objects from web pages is very
important to support effective cross-modality search.  For instance, given
an image of a painting of Monet, it is hard, if not impossible, for a
content-based system to extract the information such as ``this painting
was drawn by Monet,'' and ``it is an art of impressionism.''  However, by
analyzing the surrounding text and anchor text on links to this page, we
are able to compute this meta data.

We study how to construct the meta data of a multimedia object effectively
by using its surrounding text and anchor text of its links.  The following
are some challenges and approaches we will take:

(1) Extracting representative keywords in surrounding text.  In what range
    should we extract text that might provide representative keywords for
    a given object?  How to extract these keywords?  We need to compute a
    ``feature'' of the surrounding texts of an object, such as the range
    of the text, the sequence of each word, etc.  We need to choose the
    feature carefully, and then choose the important keywords
    correspondingly.

(2) Computing meta data using common patterns.  Similar to Task 1, we can
    use the fact that the ways multimedia objects are placed in web pages
    share some common layouts.  For instance, often people put a caption
    below an image to describe the content of the image.  If we can
    extract the common patterns these objects are placed, we can use these
    patterns to help us extract useful meta data for other objects.  That
    is, we start from some seed objects, together with their meta data,
    and see how they appear in web pages.  Then we compute the patterns of
    these objects, and apply these patterns to web pages to extract more
    objects with their meta data.  We iteratively extract patterns and
    objects (with their meta data).

    This approach is similar to the idea in \cite{Brin}.  However, the
    following challenges make our problem unique:

\begin{itemize}
   \item How to represent the layout of an object in a page?  To what
   granularity do we want to describe the layout?  This problem also
   exists in Task 1.  In \cite{Brin}, we can use a regular expression to
   represent an entity (e.g., a book title) is mentioned in a sentence.
   Here we are dealing with multimedia objects at a higher level.

   \item How to identify a multimedia object consistently in different
   pages?  For example, the same image could be in different formats in
   different pages, such as JPG, GIF, and BMP.  Some images could have
   transformed slightly.  Our approach should be able to identify the
   multiple instances of the same object in different pages, even if they
   have different binary codes.
\end{itemize}

(3) Mining anchor text on links to construct meta data.  Anchor text on
    links to a multimedia object provides very rich content about the
    object.  In addition, different web pages have different ``authority''
    in identifying the content of an object.  So we can utility the
    results of Hub/Authority or PageRank to weight the importance of the
    anchor text from different pages.  By using a weighting function on
    different links to an object, we can compute the meta data accurately.

\subsubsection*{Task 3: Computing cross-modality result using ranking functions}

To do a cross-modality query, we often need to combine information of text
and multimedia objects.  Thus it is critical to have an effective ranking
function for these different types of information.  For instance, suppose
a user wants to find all the web pages that contain a Monet painting image
and are related to impressionism.  To answer this query, we need to access
text information to find pages related to impressionism, and access image
information to find images of Monet paintings.  We should choose the right
weighting function for these different types of information to compute the
right relevance of each page to a query.

In addition, given a ranking function, how to compute the results very
efficiently is also a thrust of our research.  One challenge is that the
data size is large.  Currently there are more than 1.6 billion pages on
yhr Web.  Furthermore, the number of pages is growing very rapidly.  Our
searching algorithm should be efficient and scalable.  In the case where
it takes too much time to compute the answer, the algorithm will return
some approximate answers without compromising the accuracy too much.

The algorithm can use user feedback to improve its efficiency and
accuracy.  The user could either explicitly specify what pages are good or
bad, or click on pages, or the system will automatically detect what pages
are more preferred by the user.  Then the algorithm will adjust its
weighting function based on the user feedback to improve the result
quality for future queries.  In this sense, the algorithm is {\\em
personalized}, since it automatically adjusts its weighting function based
on the feedback from different users.

In addition, to support the weighting function, we need to design
efficient indexing structures for different modalities.

\subsubsection*{Related work}

{\em Data Integration}: The rapid growth of the Internet is giving us
access to an unprecedented number of heterogeneous information sources.
Many information-integration systems (e.g., TSIMMIS \cite{Chawathe94-1},
the Information Manifold \cite{Levy96-1}, Garlic \cite{Roth97-1},
Infomaster \cite{Genesereth97-1}, Disco \cite{Tomasic98-1}, Tukwila
\cite{Ives99-1}, and InfoSleuth \cite{Bayardo97-1}) have been proposed to
support seamless access to these heterogeneous data sources.  These
systems build {\em wrappers} on web sources to extract information from
web pages.

{\em Multimedia Ranking}: Several techniques (\cite{Fagin,Gravano,Chang})
have been proposed on how to rank information from different sources.


%Related work: web-page extraction.  Wrapper generation.  Chen's work on
%integration.
